\documentclass[12pt, titlepage]{article}

\usepackage{amsmath, mathtools}

\usepackage[round]{natbib}
\usepackage{amsfonts}
\usepackage{amssymb}
\usepackage{graphicx}
\usepackage{colortbl}
\usepackage{xr}
\usepackage{hyperref}
\usepackage{longtable}
\usepackage{xfrac}
\usepackage{tabularx}
\usepackage{float}
\usepackage{siunitx}
\usepackage{booktabs}
\usepackage{multirow}
\usepackage[section]{placeins}
\usepackage{caption}
\usepackage{fullpage}

\hypersetup{
bookmarks=true,     % show bookmarks bar?
colorlinks=true,       % false: boxed links; true: colored links
linkcolor=red,          % color of internal links (change box color with linkbordercolor)
citecolor=blue,      % color of links to bibliography
filecolor=magenta,  % color of file links
urlcolor=cyan          % color of external links
}

\newcommand{\srsref}[2]{\href{https://github.com/Nour-Fawaz/MacaBoca/blob/main/pdfs/SRS/SRS.pdf}{#1\ref*{#2}}}
\newcommand{\hazref}[1]{\href{https://github.com/Nour-Fawaz/MacaBoca/blob/main/pdfs/HazardAnalysis/HazardAnalysis.pdf}{#1}}

\usepackage{array}

\externaldocument{../../SRS/SRS}

\input{../../Comments.text}
\input{../../Common.text}

\begin{document}

\title{Module Interface Specification for \progname{}}

\author{\authname}

\date{\today}

\maketitle

\pagenumbering{roman}

\section{Revision History}

\begin{tabularx}{\textwidth}{p{3cm}p{2cm}X}
\toprule {\bf Date} & {\bf Version} & {\bf Notes}\\
\midrule
Date 1 & 1.0 & Notes\\
Date 2 & 1.1 & Notes\\
\bottomrule
\end{tabularx}

~\newpage

\section{Symbols, Abbreviations and Acronyms}

See SRS Documentation at \url{https://github.com/Nour-Fawaz/MacaBoca/blob/main/pdfs/SRS/SRS.pdf}

\wss{Also add any additional symbols, abbreviations or acronyms}

\newpage

\tableofcontents

\newpage

\pagenumbering{arabic}

\section{Introduction}

The following document details the Module Interface Specifications for
\wss{Fill in your project name and description}

Complementary documents include the System Requirement Specifications
and Module Guide.  The full documentation and implementation can be
found at \url{...}.  \wss{provide the url for your repo}

\section{Notation}

\wss{You should describe your notation.  You can use what is below as
  a starting point.}

The structure of the MIS for modules comes from \citet{HoffmanAndStrooper1995},
with the addition that template modules have been adapted from
\cite{GhezziEtAl2003}.  The mathematical notation comes from Chapter 3 of
\citet{HoffmanAndStrooper1995}.  For instance, the symbol := is used for a
multiple assignment statement and conditional rules follow the form $(c_1
\Rightarrow r_1 | c_2 \Rightarrow r_2 | ... | c_n \Rightarrow r_n )$.

The following table summarizes the primitive data types used by \progname. 

\begin{center}
\renewcommand{\arraystretch}{1.2}
\noindent 
\begin{tabular}{l l p{7.5cm}} 
\toprule 
\textbf{Data Type} & \textbf{Notation} & \textbf{Description}\\ 
\midrule
character & char & a single symbol or digit\\
integer & $\mathbb{Z}$ & a number without a fractional component in (-$\infty$, $\infty$) \\
natural number & $\mathbb{N}$ & a number without a fractional component in [1, $\infty$) \\
real & $\mathbb{R}$ & any number in (-$\infty$, $\infty$)\\
\bottomrule
\end{tabular} 
\end{center}

\noindent
The specification of \progname \ uses some derived data types: sequences, strings, and
tuples. Sequences are lists filled with elements of the same data type. Strings
are sequences of characters. Tuples contain a list of values, potentially of
different types. In addition, \progname \ uses functions, which
are defined by the data types of their inputs and outputs. Local functions are
described by giving their type signature followed by their specification.

\section{Module Decomposition} \label{Module_Decomp}

The following table is taken directly from the Module Guide document for this project.

\begin{table}[h!]
\centering
\begin{tabular}{p{0.3\textwidth} p{0.6\textwidth}}
\toprule
\textbf{Level 1} & \textbf{Level 2}\\
\midrule

{Hardware-Hiding} & ~ \\
\midrule

\multirow{7}{0.3\textwidth}{Behaviour-Hiding} & Input Parameters\\
& Output Format\\
& Output Verification\\
& Temperature ODEs\\
& Energy Equations\\ 
& Control Module\\
& Specification Parameters Module\\
\midrule

\multirow{3}{0.3\textwidth}{Software Decision} & {Sequence Data Structure}\\
& ODE Solver\\
& Plotting\\
\bottomrule

\end{tabular}
\caption{Module Hierarchy}
\label{TblMH}
\end{table}

\newpage
~\newpage

\section{MIS of \wss{Module Name}} \label{Module} \wss{Use labels for
  cross-referencing}

\wss{You can reference SRS labels, such as R\ref{R_Inputs}.}

\wss{It is also possible to use \LaTeX for hypperlinks to external documents.}

\subsection{Module}

\wss{Short name for the module}

\subsection{Uses}


\subsection{Syntax}

\subsubsection{Exported Constants}

\subsubsection{Exported Access Programs}

\begin{center}
\begin{tabular}{p{2cm} p{4cm} p{4cm} p{2cm}}
\hline
\textbf{Name} & \textbf{In} & \textbf{Out} & \textbf{Exceptions} \\
\hline
\wss{accessProg} & - & - & - \\
\hline
\end{tabular}
\end{center}

\subsection{Semantics}

\subsubsection{State Variables}

\wss{Not all modules will have state variables.  State variables give the module
  a memory.}

\subsubsection{Environment Variables}

\wss{This section is not necessary for all modules.  Its purpose is to capture
  when the module has external interaction with the environment, such as for a
  device driver, screen interface, keyboard, file, etc.}

\subsubsection{Assumptions}

\wss{Try to minimize assumptions and anticipate programmer errors via
  exceptions, but for practical purposes assumptions are sometimes appropriate.}

\subsubsection{Access Routine Semantics}

\noindent \wss{accessProg}():
\begin{itemize}
\item transition: \wss{if appropriate} 
\item output: \wss{if appropriate} 
\item exception: \wss{if appropriate} 
\end{itemize}

\wss{A module without environment variables or state variables is unlikely to
  have a state transition.  In this case a state transition can only occur if
  the module is changing the state of another module.}

\wss{Modules rarely have both a transition and an output.  In most cases you
  will have one or the other.}

\subsubsection{Local Functions}

\wss{As appropriate} \wss{These functions are for the purpose of specification.
  They are not necessarily something that is going to be implemented
  explicitly.  Even if they are implemented, they are not exported; they only
  have local scope.}



\section{MIS of ML Transcription Model} \label{MLTransModel}
The ML Transcription Model module converts an assessment audio file into timed text segments using a locally run Whisper model, flagging uncertain segments for clinician review (\srsref{SR}{SR_ClinicianFinalJudgment}, \srsref{NFR}{NFR-S3}, \srsref{NFR}{NFR-A1}).
\subsection{Module}
MLTransModel

\subsection{Uses}
\begin{itemize}
\item Local file system (to read the audio file)
\item Locally installed Whisper speech-to-text model
\end{itemize}
This module is used by the SAFA module, which supplies the audio file and displays the returned transcript. Storage and transmission of audio and transcripts are handled by other modules (Database, Encryption).

\subsection{Syntax}

\subsubsection{Exported Constants}
\begin{itemize}
\item CHUNK\_DURATION: $\mathbb{R}$\\
Length (in seconds) of each increment of audio fed to the model. Each increment yields one segment.
\item CONFIDENCE\_THRESHOLD: $\mathbb{R}$\\
Segments with confidence below this value are flagged for review.
\item SUPPORTED\_AUDIO\_FORMATS: set of String\\
Audio file formats the module accepts.
\end{itemize}
This module also uses the auxiliary constants MIN AUDIO LEVEL ($-30$ dBFS average) and MAX REC LEN (30 minutes), defined in the SRS.

\subsubsection{Exported Access Programs}
\begin{center}
\begin{tabular}{p{3cm} p{3.5cm} p{4cm} p{4cm}}
\hline
\textbf{Name} & \textbf{In} & \textbf{Out} & \textbf{Exceptions} \\
\hline
transcribe & audioPath: String & transcript: sequence of Segment & InvalidAudioError, ModelUnavailableError \\
\hline
\end{tabular}
\end{center}

\noindent The access programs use the following type:
\begin{itemize}
\item Segment: a tuple (text: String, start: $\mathbb{R}$, end: $\mathbb{R}$, confidence: $\mathbb{R}$, flagged: $\mathbb{B}$), representing one displayable chunk of transcribed speech.
  \begin{itemize}
  \item text is the transcribed speech.
  \item start and end give the time range of the chunk in the recording.
  \item confidence is a score in $[0,1]$.
  \item flagged is true iff the segment needs clinician review.
  \end{itemize}
\item The transcript is a sequence of Segment, ordered by start time.
\end{itemize}

\subsection{Semantics}

\subsubsection{State Variables}
None. The module retains no audio or transcript between calls.

\subsubsection{Environment Variables}
\begin{itemize}
\item audioFile: the audio file on the local file system that is being transcribed.
\item whisperModel: the locally installed speech-to-text model that converts a chunk of audio into text with a confidence score.
\end{itemize}

\subsubsection{Assumptions}
\begin{itemize}
\item The model is installed locally and available before transcribe is called.
\item audioPath refers to an audio file of one complete assessment session, supplied by the SAFA module.
\item The model reports a confidence score for each chunk it transcribes.
\item The SAFA module presents flagged segments to the clinician, who reviews them against the video playback before the transcript is used in an assessment (\hazref{HA-MLTM1}).
\item This module does not transmit audio or transcripts; transcription runs entirely on the local device, and storage and transmission are handled by the Database and Encryption modules (\hazref{HA-MLTM2}, NFR\ref{NFR-S3}).
\item The audio file is no longer than MAX\_REC\_LEN.
\end{itemize}

\subsubsection{Access Routine Semantics}

\noindent transcribe(audioPath):
\begin{itemize}
\item transition: None
\item output: A sequence of Segments.
  \begin{itemize}
  \item Divide the audio at audioPath into consecutive chunks of CHUNK\_DURATION (the last may be shorter).
  \item For each chunk in order, obtain its text and confidence from whisperModel and form a Segment with the chunk's time range. Set flagged to the value of \texttt{shouldFlag} for that segment and chunk.
  \item The output is these Segment(s) in time order.
  \end{itemize}
\item exception:
  \begin{itemize}
  \item InvalidAudioError if the file does not exist, is unreadable, is empty, or its format is not in SUPPORTED\_AUDIO\_FORMATS.
  \item ModelUnavailableError if whisperModel cannot be loaded or run.
  \end{itemize}
\end{itemize}

\subsubsection{Local Functions}

\noindent shouldFlag(s: Segment, chunkAudio):
\begin{itemize}
\item transition: None
\item output: true iff $s.confidence <$ CONFIDENCE\_THRESHOLD, or the average audio level of chunkAudio is below MIN\_AUDIO\_LEVEL (the model cannot reliably recognize speech at that level, so text may be invented or speech missed).
\item exception: None
\end{itemize}

\newpage

\section{MIS of Video Recording Module} \label{VideoRecording}
The Video Recording module captures an assessment session from the patient's camera and microphone, checks recording quality beforehand, and ensures the recording is saved intact and playable (\srsref{NFR}{NFR-M1}, \srsref{NFR}{NFR-P1}).

\subsection{Module}
VideoRecording

\subsection{Uses}
\begin{itemize}
\item Web Platform module (camera and microphone access, and temporary storage on the patient's device)
\item Video Storage module (permanent storage, retrieval and backup of recordings)
\end{itemize}
This module is used by the Web Platform module during an assessment. Encryption of stored recordings is handled by the Encryption module.

\subsection{Syntax}

\subsubsection{Exported Constants}
\begin{itemize}
\item MIN\_BRIGHTNESS: $\mathbb{R}$\\
Minimum average video brightness for the lighting to be considered adequate.
\item MIN\_FRAME\_RATE: $\mathbb{R}$\\
Minimum video frame rate (frames per second) accepted.
\item MAX\_UPLOAD\_ATTEMPTS: $\mathbb{N}$\\
Number of times saving is attempted automatically before the failure is reported.
\end{itemize}
This module also uses the auxiliary constants MIN\_AUDIO\_LEVEL ($-30$ dBFS average), MIN\_VIDEO\_RESOLUTION ($640 \times 480$ pixels) and MAX\_REC\_LEN (30 minutes), defined in the SRS.

\subsubsection{Exported Access Programs}
\begin{center}
\begin{tabular}{p{3cm} p{3.5cm} p{4cm} p{4.5cm}}
\hline
\textbf{Name} & \textbf{In} & \textbf{Out} & \textbf{Exceptions} \\
\hline
checkQuality & - & report: QualityReport & PermissionDeniedError, DeviceUnavailableError \\
startRecording & - & - & PermissionDeniedError, DeviceUnavailableError, RecordingInProgressError \\
stopRecording & - & - & NotRecordingError \\
saveRecording & assessmentID: String & - & NoRecordingError, UploadFailedError, VerificationFailedError \\
getStatus & - & status: \{IDLE, RECORDING, PENDING\_SAVE, SAVED\} & - \\
\hline
\end{tabular}
\end{center}

\noindent The access programs use the following type:
\begin{itemize}
\item QualityReport: a tuple (passed: $\mathbb{B}$, issues: sequence of String).
  \begin{itemize}
  \item passed is true iff the audio and video meet the minimum levels above.
  \item issues describes each problem found (e.g., audio too quiet, lighting too dark, resolution too low) so the patient can correct their setup.
  \end{itemize}
\end{itemize}
\subsection{Semantics}

\subsubsection{State Variables}
\begin{itemize}
\item status: \{IDLE, RECORDING, PENDING\_SAVE, SAVED\}\\
Where the current recording is in its lifecycle.
\item localCopy: Recording or none\\
The temporary copy of the current recording held on the patient's device until the system confirms it was saved successfully.
\end{itemize}

\subsubsection{Environment Variables}
\begin{itemize}
\item camera: the patient's video input device.
\item microphone: the patient's audio input device.
\item deviceStorage: temporary storage on the patient's device, in which localCopy is held.
\end{itemize}

\subsubsection{Assumptions}
\begin{itemize}
\item The caller runs checkQuality, and the patient has corrected any reported issues, before startRecording is called (\hazref{HA-VR1}).
\item The Video Storage module reports success or failure of a store request, can return a stored recording for verification, and can create a backup copy of a stored recording.
\item The patient's device has enough temporary storage to hold the full recording.
\item assessmentID is supplied by the caller and identifies the assessment the recording belongs to.
\end{itemize}

\subsubsection{Access Routine Semantics}

\noindent checkQuality():
\begin{itemize}
\item transition: None
\item output: A QualityReport based on a short sample of the camera and microphone input. passed is true iff the sample's audio level is at least MIN\_AUDIO\_LEVEL, its brightness is at least MIN\_BRIGHTNESS, and its resolution and frame rate are at least MIN\_RESOLUTION and MIN\_FRAME\_RATE. Each failed condition adds an entry to issues. The sample is discarded.
\item exception: PermissionDeniedError if access to the camera or microphone is not granted; DeviceUnavailableError if either device is missing or in use elsewhere.
\end{itemize}

\noindent startRecording():
\begin{itemize}
\item transition: status := RECORDING; begin capturing audio and video from camera and microphone.
\item output: None
\item exception: PermissionDeniedError if access to the camera or microphone is not granted; DeviceUnavailableError if either device is missing or in use elsewhere; RecordingInProgressError if status = RECORDING.
\end{itemize}

\noindent stopRecording():
\begin{itemize}
\item transition: Stop capturing. localCopy := the complete recording, saved in deviceStorage. status := PENDING\_SAVE.
\item output: None
\item exception: NotRecordingError if status $\neq$ RECORDING.
\end{itemize}

\noindent saveRecording(assessmentID):
\begin{itemize}
\item transition:
  \begin{itemize}
  \item Request that Video Storage store localCopy under assessmentID, trying up to MAX\_UPLOAD\_ATTEMPTS times.
  \item Retrieve the stored recording and check it with \texttt{isIntact}.
  \item If the check passes, request that Video Storage create a backup of the stored recording, then localCopy := none and status := SAVED.
  \item If any step fails, localCopy and status are unchanged, so saveRecording can be called again.
  \end{itemize}
\item output: None
\item exception: NoRecordingError if status $\neq$ PENDING\_SAVE; UploadFailedError if storing the recording did not succeed after MAX\_UPLOAD\_ATTEMPTS attempts; VerificationFailedError if the stored recording is incomplete or cannot be played back.
\end{itemize}

\noindent getStatus():
\begin{itemize}
\item transition: None
\item output: status
\item exception: None
\end{itemize}

\subsubsection{Local Functions}

\noindent isIntact(stored: Recording, original: Recording):
\begin{itemize}
\item transition: None
\item output: true iff stored can be decoded and played from start to end, and its duration and content match original.
\item exception: None
\end{itemize}

\newpage
\section{MIS of Transcription Service Interface Module} \label{TranscriptionServiceInterface}
The Transcription Service Interface module generates a draft transcript for an assessment recording and manages the clinician's review of it: edits, speaker roles and verification (\srsref{FR}{FR_AutomaticTranscription}, \srsref{FR}{FR_TranscriptionEditing}, \srsref{FR}{FR_Label6}, \srsref{NFR}{NFR-RE1}).

\subsection{Module}
TranscriptionServiceInterface

\subsection{Uses}
\begin{itemize}
\item MLTransModel module (generates timed text segments from audio)
\item Video Storage module (provides the recording of an assessment)
\item Database module (stores transcripts and records changes to them in the audit log)
\end{itemize}
This module is used by the Web Platform module, which displays the transcript and passes the clinician's actions to it.

\subsection{Syntax}

\subsubsection{Exported Constants}
\begin{itemize}
\item SPEAKER\_ROLES: set of String = \{System, Child, Parent, Experimenter, Multiple Speakers, Others\}\\
The roles a clinician can assign to a speaker.
\item GENERATION\_ATTEMPTS: $\mathbb{N}$\\
Number of times transcription is attempted automatically before a failure is reported, so the clinician does not have to restart it manually.
\end{itemize}
This module also relies on the auxiliary constants MAX\_REC\_LEN (30 minutes) and SYNC\_TOL (500 ms), defined in the SRS.

\subsubsection{Exported Access Programs}
\begin{center}
\begin{tabular}{p{3.2cm} p{4.3cm} p{3cm} p{4.5cm}}
\hline
\textbf{Name} & \textbf{In} & \textbf{Out} & \textbf{Exceptions} \\
\hline
generateTranscript & assessmentID: String, userID: String, portion: Interval or none & - & AssessmentNotFoundError, RecordingUnavailableError, TranscriptionInProgressError, PortionAlreadyTranscribedError, GenerationFailedError \\
getTranscript & assessmentID: String, userID: String & transcript: Transcript & AssessmentNotFoundError, NoTranscriptError \\
editSegmentText & assessmentID: String, segmentID: String, newText: String, userID: String & - & NoTranscriptError, SegmentNotFoundError, TranscriptionInProgressError, SaveFailedError \\
assignSpeaker & assessmentID: String, segmentID: String, role: String, userID: String & - & NoTranscriptError, SegmentNotFoundError, InvalidRoleError, TranscriptionInProgressError, SaveFailedError \\
verifySegment & assessmentID: String, segmentID: String, userID: String & - & NoTranscriptError, SegmentNotFoundError, TranscriptionInProgressError, SaveFailedError \\
completeTranscript & assessmentID: String, userID: String & unverified: sequence of String & NoTranscriptError, TranscriptionInProgressError, SaveFailedError \\
\hline
\end{tabular}
\end{center}
\noindent The access programs use the following types:
\begin{itemize}
\item Interval: a pair of times (from, to) in seconds on the recording's timeline.
\item Transcript: a pair (status, segments).
  \begin{itemize}
  \item status is one of \{BEING\_GENERATED, DRAFT, VERIFIED\}.
  \item segments is a sequence of Segment, ordered by start time.
  \end{itemize}
\item Segment: a tuple (segmentID: String, start: $\mathbb{R}$, end: $\mathbb{R}$, text: String, originalText: String, speaker: String or none, language: String, confidence: $\mathbb{R}$, flagged: $\mathbb{B}$, verified: $\mathbb{B}$).
  \begin{itemize}
  \item Times are stored to the nearest tenth of a second or better so playback can be synchronized within SYNC\_TOL (\srsref{NFR}{NFR-P2}).
  \item A segment counts as edited when text differs from originalText, so automatic and clinician-modified content can be told apart (\srsref{NFR}{NFR-UN1}).
  \end{itemize}
\end{itemize}
\subsection{Semantics}

\subsubsection{State Variables}
None. All transcript data is held by the Database module, so work saved in one session is available in the next (\srsref{NFR}{NFR-RE2}).

\subsubsection{Environment Variables}
None.

\subsubsection{Assumptions}
\begin{itemize}
\item The caller has confirmed that userID belongs to the clinician assigned to the patient of the assessment (\srsref{NFR}{NFR-S1}).
\item Each assessment has one recording, up to MAX\_REC\_LEN long, retrievable from Video Storage.
\item The Database module records an audit entry for every view or change of a transcript made through this module (\srsref{NFR}{NFR-S4}).
\item Clinicians review flagged segments against the video playback before relying on the transcript (\hazref{HAMLTM1}).
\item Encryption of recordings and transcripts in transit and in storage is provided by the Encryption, Video Storage and Database modules (\srsref{NFR}{NFR-S3}).
\end{itemize}

\subsubsection{Access Routine Semantics}

\noindent generateTranscript(assessmentID, userID, portion):
\begin{itemize}
\item transition:
  \begin{itemize}
  \item Set the transcript status to BEING\_GENERATED.
  \item Obtain the assessment's recording from Video Storage and extract the audio of the whole recording, or of portion if given, into a temporary audio file (\texttt{extractAudio}).
  \item Call MLTransModel.transcribe on the audio file, retrying up to GENERATION\_ATTEMPTS times if the model is unavailable.
  \item Store each returned segment with text and originalText equal to the generated text, speaker none, verified false, and its times shifted by the start of portion (no shift for the whole recording).
  \item Set the transcript status to DRAFT.
  \item Delete the temporary audio file, whether or not generation succeeded.
  \item If generation fails, remove any segments added by this call and restore the previous status.
  \end{itemize}
\item output: None
\item exception: AssessmentNotFoundError if assessmentID does not exist; RecordingUnavailableError if the recording cannot be obtained; TranscriptionInProgressError if the status is already BEING\_GENERATED; PortionAlreadyTranscribedError if existing segments overlap portion (or any segments exist, when portion is none); GenerationFailedError if all attempts fail.
\end{itemize}

\noindent getTranscript(assessmentID, userID):
\begin{itemize}
\item transition: None
\item output: The stored transcript, including its status (so the caller can tell the clinician that it is still being generated).
\item exception: AssessmentNotFoundError if assessmentID does not exist; NoTranscriptError if no transcript has been generated for it.
\end{itemize}

\noindent editSegmentText(assessmentID, segmentID, newText, userID):
\begin{itemize}
\item transition: Set text of the segment to newText. originalText is unchanged. The change is saved immediately.
\item output: None
\item exception: NoTranscriptError, SegmentNotFoundError, TranscriptionInProgressError if the status is BEING\_GENERATED; SaveFailedError if the change could not be saved, in which case nothing is changed and the caller keeps the clinician's input to retry.
\end{itemize}

\noindent assignSpeaker(assessmentID, segmentID, role, userID):
\begin{itemize}
\item transition: Set speaker of the segment to role. The change is saved immediately.
\item output: None
\item exception: InvalidRoleError if role is not in SPEAKER\_ROLES; otherwise as for editSegmentText.
\end{itemize}

\noindent verifySegment(assessmentID, segmentID, userID):
\begin{itemize}
\item transition: Set verified of the segment to true. The change is saved immediately.
\item output: None
\item exception: As for editSegmentText.
\end{itemize}

\noindent completeTranscript(assessmentID, userID):
\begin{itemize}
\item transition: Set the transcript status to VERIFIED. Segments that are not verified keep verified false.
\item output: The segmentIDs of the segments that are still unverified, so the caller can offer to return to them.
\item exception: NoTranscriptError, TranscriptionInProgressError, SaveFailedError.
\end{itemize}

\subsubsection{Local Functions}

\noindent extractAudio(recording, portion):
\begin{itemize}
\item transition: Writes a temporary audio file.
\item output: A file in a format supported by MLTransModel containing the audio of recording (restricted to portion, if given).
\item exception: RecordingUnavailableError if the recording has no readable audio.
\end{itemize}

\newpage

\bibliographystyle {plainnat}
\bibliography {../../../refs/References}

\newpage

\newpage{}

\section{Appendix --- Generative AI Use Disclosure}
This section documents all the uses of generative artificial intelligence used throughout the document.

\subsection{AI Tools Used}

\begin{itemize}
  \item Claude Sonnet 5 (Medium)
  \item ChatGPT-5.6 Luna
\end{itemize}

\subsection{How and Where Used}

The following lists the uses of AI applied throughout the entire document:
\begin{itemize}
  \item Fixing grammar and spelling
  \item Debugging possible \LaTeX{} errors
  \item Performing a final check over the document at the end, comparing each section against the original document comments and the rubric to ensure nothing was missed
  \item Removing template comments
\end{itemize}

\noindent\textbf{Section~\ref{Module_Decomp}}
\begin{itemize}
    \item Format the table of existing modules for \LaTeX{}
\end{itemize}

\noindent\textbf{All Modules}
\begin{itemize}
    \item Brainstorm modules in addition to the already existing ones
    \item Identify possible state and environment variables for each module
\end{itemize}

\subsection{Verification of AI-Generated Content}

To validate the LLM's output, each team member reviewed their assigned section and presented it to the group, summarizing the work they had completed since the last meeting.
All members were able to explain their sections clearly.
Before any pull request was approved, a team member who had not authored the section reviewed it to confirm that it was cohesive, accurate, and consistent with the rest of the document.
Because the document contained no factual claims, the LLM's output did not require any additional fact-checking.

\newpage{}

\section{Appendix --- Reflection}

\wss{Not required for CAS 741 projects}

The information in this section will be used to evaluate the team members on the
graduate attribute of Problem Analysis and Design.

\input{../../Reflection.text}

\begin{enumerate}
  \item What went well while writing this deliverable? 
  \item What pain points did you experience during this deliverable, and how
    did you resolve them?
  \item Which of your design decisions stemmed from speaking to your client(s)
  or a proxy (e.g. your peers, stakeholders, potential users)? For those that
  were not, why, and where did they come from?
  \item While creating the design doc, what parts of your other documents (e.g.
  requirements, hazard analysis, etc), it any, needed to be changed, and why?
  \item What are the limitations of your solution?  Put another way, given
  unlimited resources, what could you do to make the project better? (LO\_ProbSolutions)
  \item Give a brief overview of other design solutions you considered.  What
  are the benefits and tradeoffs of those other designs compared with the chosen
  design?  From all the potential options, why did you select the documented design?
  (LO\_Explores)
  \item (After you have implemented another team's module, which means this
  isn't filled in until after the original deadline). What did you learn by
  implementing another team's module? Were all the details you needed in the
  documentation, or did you need to make assumptions, or ask the other team
  questions?  If your team also had another team implement one of your modules,
  what was this experience like?  Are there things in your documentation you
  could have changed to make the process go more smoothly for when an
  ``outsider'' completes some of the implementation?

  Your reflection should be about concepts learned in your SE education, such as modular decomposition, cohesion, coupling, information hiding, etc.  You should critique the quality and utility of the specification you received.

  You should create a PR for your implemented module.  A link to the PR should be part of your reflection.

\end{enumerate}

\end{document}